\section{The trust layer}
\label{sec:trust}

The trust layer reads one \nts{} per sentence attempt and returns one action. It never reads raw signals or decoder internals. \Cref{fig:decision} shows the decision flow.

\subsection{Calibration}
Decoder log-probabilities are often over-confident. We fit a frame temperature $T_f$ that rescales every frame, $\tilde p_t(v) \propto p_t(v)^{1/T_f}$, by minimising the negative log-likelihood of the truth on fit data, and a sentence temperature $T_s$ for the sentence-level confidence $\hat p_{\mathrm{top1}}$, the normalised probability of the top candidate among the n-best \citep{guo2017calibration}. We report the expected calibration error (ECE) of $\hat p_{\mathrm{top1}}$ before and after.

\subsection{Candidates and neural-evidence guard}
Candidate sentences come from a CTC prefix beam search with a character 6-gram LM and a word unigram, both trained on Tatoeba sentences, and an n-best list (n = 10) \citep{graves2006ctc,hannun2014deepspeech}. For the published intracortical sentences, the candidate is the given text. Because the LM can write fluent text from any input, a guard scores how much of each candidate the LM invented. It combines two signals: the share of the candidate's characters whose posterior under the calibrated frames alone (no LM) is low (\emph{low evidence}), and one minus the posterior-weighted word agreement between the top candidate and the other n-best candidates (\emph{n-best inconsistency}); the guard risk is the maximum of the two weighted signals. A candidate with high guard risk cannot be committed.

\subsection{Decision policy}
The policy chooses one of four actions: \act{commit} the top candidate; \act{offer} the top $k$ candidates for the user to pick; \act{offer\_words}, an offer with uncertain words highlighted; or \act{abstain}, after which the user types the sentence with the fallback keyboard. The thresholds (commit probability, guard risk, offer size, abstain risk, fallback) are chosen on fit data from a small grid to maximise \ewpm{} (\cref{sec:protocol}) subject to a wrong-commit rate $\le 0.05$ on the fit data. This plug-in fit has no guarantee on new data, which motivates \cref{sec:crc}.

\subsection{Finite-sample risk control for the commit threshold}
\label{sec:crc}
Let $s_i$ be the score of item $i$ ($\hat p_{\mathrm{top1}}$ if the other policy rules allow a commit, else $-\infty$) and $\lambda$ the commit threshold. The loss is $L_i(\lambda) = c_i$ if $s_i \ge \lambda$ (with $c_i = 1$ if the top candidate is wrong) and $f_i$ otherwise (the loss of the non-commit action, 0 here), made monotone in $\lambda$ by taking $\max(c_i, f_i)$. On $n$ calibration items with known truth, conformal risk control \citep{angelopoulos2024crc} chooses
\begin{equation}
  \hat\lambda_{\mathrm{CRC}} = \inf\Bigl\{\lambda : \tfrac{n}{n+1}\hat R_n(\lambda) + \tfrac{1}{n+1} \le \alpha\Bigr\},
\end{equation}
where $\hat R_n$ is the empirical risk. If calibration and test items are exchangeable, $\mathbb{E}[L_{n+1}(\hat\lambda_{\mathrm{CRC}})] \le \alpha$. No feasible $\lambda$ exists for $n < 1/\alpha - 1$ (19 for $\alpha = 0.05$); the layer then never commits. Learn-then-Test \citep{angelopoulos2021ltt} instead tests $H_0\colon R(\lambda) > \alpha$ along a decreasing grid with exact binomial p-values and fixed-sequence testing, which gives $\Pr(R(\hat\lambda_{\mathrm{LTT}}) \le \alpha) \ge 1 - \delta$. Any threshold above $\hat\lambda$ keeps the guarantee, so the layer can combine CRC with the \ewpm{}-tuned threshold as $\hat\lambda = \max(\lambda_{\mathrm{tuned}}, \hat\lambda_{\mathrm{CRC}})$. The other parts of the pipeline (temperatures, guard, other thresholds) are fitted on separate data. We use $\alpha = 0.05$ and $\delta = 0.1$ throughout.

\subsection{Online re-calibration under drift}
\label{sec:online}
Exchangeability fails when a session drifts. The online controller re-fits the threshold during the session from user feedback (\cref{fig:online}). Which feedback exists depends on the action: a corrected commit reveals a wrong commit; an uncorrected commit is right only if silence means right; an offer reveals the truth when the user picks or retypes; an abstention reveals nothing unless the final text is logged. The controller therefore accepts two feedback levels: the \emph{realised loss} of committed items, and, when the final text is known, the counterfactual commit loss of non-committed items. Feedback can arrive late and out of order. \Cref{tab:online} lists the four methods.

\begin{table}[t]
  \centering
  \caption{\textbf{Online re-calibration methods for the commit threshold.} $L_t$: realised loss at item $t$; $T$: number of items in the session. The two \texttt{aci} variants bound the long-run average loss for any sequence of items, without exchangeability.}
  \label{tab:online}
  \small
  \begin{adjustbox}{max width=\linewidth}\begin{tabular}{@{}llll@{}}
    \toprule
    \textbf{Method} & \textbf{Update} & \textbf{Guarantee} & \textbf{Feedback} \\
    \midrule
    \texttt{window} & CRC on the last $N$ items with known truth & approximate & full \\
    \texttt{decay} & weighted CRC, weights $\rho^{\text{age}}$ \citep{barber2023beyond,farinhas2024nonexch} & approximate & full \\
    \texttt{aci} & $\lambda_{t+1} = \max(\lambda_{\min}, \lambda_t + \eta(L_t - \alpha))$ & $\frac1T\sum L_t \le \alpha + \frac{\lambda_{T+1} - \lambda_1}{\eta T}$ & realised only \\
    \texttt{aci\_window} & window CRC at level $\alpha_t$; $\alpha_{t+1} = \alpha_t + \gamma(\alpha - L_t)$ & $\frac1T\sum L_t \le \alpha + \frac{\alpha_1 - \alpha_{T+1}}{\gamma T}$ & full + realised \\
    \bottomrule
  \end{tabular}\end{adjustbox}
\end{table}

The \texttt{aci} rule follows adaptive conformal inference and online risk control \citep{gibbs2021aci,angelopoulos2023pid,feldman2023rolling}; \texttt{aci\_window} applies the ACI update to the level of a windowed CRC \citep{gibbs2021aci}. We use $N = 200$, $\rho = 0.99$, $\eta = 0.05$ and $\gamma = 0.005$, fixed before the evaluation. These bounds are long-run averages over the $T$ items of a session; they say nothing about short windows, and they cover only the errors that the system observes. A wrong commit that the user does not notice (a \emph{silent error}) is outside them.

\subsection{Provenance ledger}
\label{sec:ledger}
Each decision is appended to a JSON Lines ledger. Entry 0 is a genesis record; \emph{context} entries are content-addressed (their id is the SHA-256 of the context) and hold the adapter, decoder (name, version, weights hash, training data and licence), LM, calibration (temperatures, risk-control method, $\alpha$, $\delta$, threshold, $n$) and policy versions; \emph{decision} entries hold the action, scores, a hash of the \nts{} and its frame range; \emph{correction} entries link to the decision they correct and store the edit distance. Each entry stores \code{prev\_hash} and its own SHA-256 over the canonical JSON, so any edit, deletion or reordering breaks the chain; a separately stored head anchor detects truncation. The ledger stores no raw signals (numeric arrays longer than 64 values are rejected), stores text as a keyed hash by default, and pseudonymises subject and session ids. It is tamper-evident, not tamper-proof: there is no signature in v0.

\begin{figure}[t]
  \centering
  \resizebox{\linewidth}{!}{% Trust decision flow for one sentence: COMMIT / OFFER / ABSTAIN with the CRC threshold.
\begin{tikzpicture}[
  font=\small,
  box/.style={draw, rounded corners=2pt, align=center, fill=white, minimum height=8mm, inner sep=3pt},
  dec/.style={draw, diamond, aspect=2.4, align=center, fill=blue!6, inner sep=1pt, font=\scriptsize},
  act/.style={draw, rounded corners=6pt, align=center, minimum width=22mm, minimum height=9mm, very thick},
  arr/.style={-{Stealth[length=2mm]}, thick},
  node distance=5mm and 8mm]
  \node[box, fill=blue!10] (in) {NTS of one sentence\\\scriptsize calibrated $\log p_t(v)$};
  \node[box, right=of in] (cand) {$K$-best candidates\\\scriptsize beam + LM; $\hat p_{\mathrm{top1}}$, guard risk};
  \node[dec, right=of cand] (g) {guard risk\\$\le r_{\max}$?};
  \node[dec, right=10mm of g] (lam) {$\hat p_{\mathrm{top1}} \ge \hat\lambda$?};
  \node[act, draw=teal!70!black, fill=teal!10, right=10mm of lam] (commit) {\textsc{commit}\\\scriptsize top-1 sentence};
  \node[dec, below=9mm of g] (off) {useful\\alternatives?};
  \node[act, draw=orange!80!black, fill=orange!10, right=10mm of off] (offer) {\textsc{offer}\\\scriptsize pick from $k$, or words};
  \node[act, draw=gray, fill=gray!10, below=6mm of offer] (abst) {\textsc{abstain}\\\scriptsize user types (gaze)};
  \node[box, fill=violet!8, above=7mm of lam, text width=44mm] (thr) {\scriptsize $\hat\lambda = \max(\lambda_{\mathrm{tuned}}, \lambda_{\mathrm{CRC}})$\\\scriptsize CRC: $\mathbb{E}[\text{wrong commit}] \le \alpha$\\\scriptsize or online $\lambda_t$ (\cref{sec:online})};
  \draw[arr] (in) -- (cand);
  \draw[arr] (cand) -- (g);
  \draw[arr] (g) -- node[above, font=\scriptsize] {yes} (lam);
  \draw[arr] (lam) -- node[above, font=\scriptsize] {yes} (commit);
  \draw[arr] (g) -- node[left, font=\scriptsize] {no} (off);
  \draw[arr] (lam) |- node[pos=0.25, right, font=\scriptsize] {no} ($(off.north)+(0,3mm)$) -- (off.north);
  \draw[arr] (off) -- node[above, font=\scriptsize] {yes} (offer);
  \draw[arr] (off) |- node[pos=0.3, left, font=\scriptsize] {no} (abst);
  \draw[arr, violet!70!black, densely dashed] (thr.south) -- (lam.north);
\end{tikzpicture}
}
  \caption{\textbf{Trust decision for one sentence.} Candidates come from the calibrated stream. A candidate is committed only if the neural-evidence guard accepts it and its calibrated confidence reaches the threshold $\hat\lambda$, which is the larger of the \ewpm{}-tuned threshold and the CRC threshold (or the online threshold $\lambda_t$). Otherwise the layer offers alternatives or abstains, and the user types with the fallback keyboard.}
  \label{fig:decision}
\end{figure}

\begin{figure}[t]
  \centering
  \resizebox{0.95\linewidth}{!}{% Online re-calibration loop of the COMMIT threshold (aci and aci_window).
\begin{tikzpicture}[
  font=\small,
  box/.style={draw, rounded corners=2pt, align=center, fill=white, minimum height=9mm, inner sep=3pt, text width=33mm},
  arr/.style={-{Stealth[length=2mm]}, thick},
  node distance=6mm and 14mm]
  \node[box, fill=violet!8] (warm) {warm start\\\scriptsize CRC on a calibration block\\\scriptsize $\lambda_1$, $\alpha_1 = \alpha$};
  \node[box, fill=blue!8, right=of warm] (dec) {decide item $t$\\\scriptsize commit iff $s_t \ge \lambda_t$};
  \node[box, fill=yellow!20, right=of dec] (user) {user acts\\\scriptsize corrects, picks or retypes};
  \node[box, fill=teal!10, below=10mm of user] (fb) {feedback\\\scriptsize realised loss $L_t$; truth if known\\\scriptsize (any order, any delay)};
  \node[box, fill=teal!10, below=10mm of dec, text width=52mm] (upd) {update\\\scriptsize \texttt{aci}: $\lambda_{t+1} = \max(\lambda_{\min}, \lambda_t + \eta (L_t - \alpha))$\\\scriptsize \texttt{aci\_window}: $\alpha_{t+1} = \alpha_t + \gamma(\alpha - L_t)$;\\\scriptsize $\lambda_{t+1}$ = window CRC at level $\alpha_{t+1}$};
  \node[box, fill=gray!8, below=10mm of warm] (cert) {certificate\\\scriptsize $\frac{1}{T}\sum_t L_t \le \alpha + O(1/T)$\\\scriptsize for any sequence};
  \draw[arr] (warm) -- (dec);
  \draw[arr] (dec) -- node[above, font=\scriptsize] {action} (user);
  \draw[arr] (user) -- (fb);
  \draw[arr] (fb) -- (upd);
  \draw[arr] (upd) -- node[right, font=\scriptsize] {$\lambda_{t+1}$} (dec);
  \draw[arr, densely dashed] (upd) -- (cert);
\end{tikzpicture}
}
  \caption{\textbf{Online re-calibration loop.} A short calibration block gives the starting threshold. After each item, the realised loss (and the truth, if the final text is known) updates the threshold. For \texttt{aci} and \texttt{aci\_window}, the running wrong-commit rate is bounded by $\alpha$ plus a term that shrinks as $1/T$, for any sequence of items.}
  \label{fig:online}
\end{figure}
